\documentclass[11pt]{article}
\usepackage[margin=1in]{geometry}
\usepackage{amsmath,amssymb,amsthm}
\usepackage[hidelinks]{hyperref}
\newtheorem{theorem}{Theorem}
\newtheorem{corollary}{Corollary}
\newcommand{\E}{\mathbb E}
\title{Inactive variables in the Gamma block-radial limit}
\author{Coordinator working extension of the retained Long--Le Blanc block limit}
\date{16 November 2026}
\begin{document}
\maketitle
\noindent\textbf{Status.}
Historical draft, superseded by the linked canonical coupled-block proof; no independent review, formal verification or
novelty claim. This extends the explicitly retained positive-exponent block
theorem in \href{../../../factorial/image-rigidity/COUPLED_BLOCK_PERIODS.tex}
{Coupled block periods}, based on the pinned Long--Le Blanc EIT VI source.
That note explicitly leaves lower terms involving extra inactive variables
outside its singleton-block corollary. No canonical proof or registry is changed.

\section{The extension}
Partition active independent positive-shape Gamma variables into nonempty
blocks $B_j$. Let $S_j$ be their sums, $A_j$ their shape sums,
$d_j\ge1$, $D=\sum_jd_j$, $C_d=\prod_jd_j^{d_j}$, and let $T^{(j)}$ be
their independent Dirichlet angular variables. Also let $Z$ be any finite
tuple of independent Gamma variables of fixed strictly positive shapes,
independent of all active blocks. Write $dT=(d_1T^{(1)},\ldots,d_rT^{(r)})$.
All polynomial coefficients below may be arbitrary complex numbers.

\begin{theorem}
For $c\ne0$, homogeneous $G$ of total degree $D-1$, and
$\deg R\le D-2$, put
\[
 F(X,Z)=c\prod_jS_j^{d_j}+G(X,Z)+R(X,Z).
\]
For a nonzero polynomial marker $H$, let $e$ be its maximum degree in
the active variables only, and let $H_{[e]}(X,Z)$ be its active-degree-$e$
part, retaining all polynomial dependence on $Z$. Then
\[
 \lim_{m\to\infty}
 \frac{\E[H(X,Z)F(X,Z)^m]}
 {m^e c^m\prod_j(A_j)_{d_jm}}
 =
 \E_{T,Z}\left[
 H_{[e]}(dT,Z)
 \exp\!\left(\frac{G(dT,0)}{cC_d}\right)\right].
 \tag{1}
\]
Zero markers are identically zero. A zero limit does not decide the
entire marked sequence.
\end{theorem}

\begin{proof}
Divide $F$ by $c$. Tilt only the active block radii by
$\prod S_j^{d_jm}$ and normalize by their Gamma moments.
The resulting independent radii have laws
$\operatorname{Gamma}(d_jm+A_j,1)$, so $S_j/m\to d_j$ in probability.
The inactive $Z$ retain their original fixed laws, so $Z/m\to0$.
The angular variables retain their Dirichlet laws.

For an active radial vector $u$ near $d$, and $v=Z/m$ near zero in
the closed positive orthant, uniformly on a compact neighborhood,
\[
 \frac{F(muT,mv)}{c\,m^D\prod u_j^{d_j}}
 =
 1+\frac{G(uT,v)}{cm\prod u_j^{d_j}}+O(m^{-2}).
\]
The $m$th powers are uniformly bounded there, converge to the
exponential in (1), and the active-scaled marker converges to
$H_{[e]}(dT,Z)$. On this neighborhood
$m^{-e}|H(muT,Z)|\le C(1+|Z|)^L$ for some fixed $L$.
The latter has finite expectation under every fixed positive-shape
Gamma law. Concentration of the tilted active radii, then dominated
convergence in $T,Z$, gives the claimed local limit.

For completeness, the complement estimate retains the inactive boundary.
Scale \emph{all} radii as $(S,Z)=m(u,v)$ and put
\[
 K_m(u,v,T)=e^{-\sum u-\sum v}
 |m^{-D}F(muT,mv)/c|.
\]
Uniformly on compact sets, including coordinate boundaries,
\[
 K_m\longrightarrow
 K_0=e^{-\sum u-\sum v}\prod u_j^{d_j}.
\]
This has its unique global maximum $M=e^{-D}C_d>0$ at
$(u,v)=(d,0)$. It need not vanish on an inactive boundary.
Outside any fixed relative neighborhood of that point, the compact
maximum is strictly smaller than $M$. The polynomial bound
\[
 K_m\le C e^{-\sum u-\sum v}(1+\sum u+\sum v)^D
\]
also gives a uniform exponential gap on the unbounded tail.
All fixed-shape radial densities are locally integrable, including
shapes below one. Polynomial markers change these estimates by only
a fixed polynomial factor and a fixed power of $m$.

If $B$ is the sum of inactive shapes, Stirling's formula gives the
scaled normalization
\[
 \frac{\prod_j\Gamma(d_jm+A_j)}
 {m^{Dm+\sum A_j+B}}
 =C_*M^m m^{-r/2-B}(1+o(1)),\qquad C_*>0.
\]
The complementary compact region and tail are therefore exponentially
negligible after normalization. This proves (1) without dividing by
the leading form away from the active neighborhood.
\end{proof}

\section{Consequences}
\begin{corollary}
Suppose the highest ordinary homogeneous part of $F$ is a single
nonconstant monomial $cX^\alpha$, with arbitrary nonnegative exponents
and $c\ne0$. Every sufficiently large Gamma moment of $F$ is nonzero,
for every fixed choice of positive Gamma shapes and arbitrary complex
strictly lower-degree coefficients.
\end{corollary}
\begin{proof}
Declare the variables with positive $\alpha_i$ active singleton blocks
and all others inactive. The unmarked limit is
\[
 \exp\!\left(
 \frac{G(\alpha,0)}{c\prod_{\alpha_i>0}\alpha_i^{\alpha_i}}
 \right)\ne0.
\]
\end{proof}
For factorial moments this includes $X^3+Q_{\le2}(X,Y,Z)$ and
$X^2Y+Q_{\le2}(X,Y,Z,W)$, with arbitrary complex lower terms.
Their leading forms can be singular and their lower terms can involve
every inactive variable.

More generally any previously proved nonzero angular block limit remains
nonzero after adjoining any number of inactive variables and arbitrary
strictly lower-degree couplings, provided the restriction $G(X,0)$ is
in the original proved nonvanishing family. This is stronger than merely
adjoining unused variables.

This proves no unrestricted FC or SFC statement, no two-moment window,
and no assertion for multiple top monomials outside the retained block
architecture. The marker $Z-1$ has zero limiting mean for a unit
exponential inactive variable, while lower couplings can give nonzero
marked moments; thus zero limits must not be promoted to zero sequences.
Finite exact identities in the companion checker are transcriptions,
not a substitute for the uniform boundary estimate above.
\end{document}
